%══════════════════════════════════════════════════════════════════════
%  CHAPTER 5 — AI Components
%══════════════════════════════════════════════════════════════════════
\chapter{AI Components}
\label{ch:ai}

This chapter details the two AI tracks embedded in stage~S3 of \padonap{}:
Track~A, an XGBoost seven-class detector, and Track~B, a Transformer+LSTM
four-horizon forecaster. Both tracks were trained on CICDDoS2019 under a
strict temporal 80/20 split; serialised models are stored in
\texttt{models\_v2/} and loaded at runtime by
\texttt{pipeline/s3\_ai/inference\_layer.py}.

%%─────────────────────────────────────────────────────────────────────
\section{Feature Engineering}
\label{sec:features}

Every 5-second window is aggregated by \texttt{pipeline/s2\_features/}
into a 17-dimensional feature vector $\mathbf{x} \in \mathbb{R}^{17}$.
Features are grouped into three families:

\begin{enumerate}
  \item \textbf{Rate features} (6): \texttt{pkt\_rate}, \texttt{byte\_rate},
        \texttt{inter\_arrival\_mean}, \texttt{inter\_arrival\_std},
        \texttt{inter\_arrival\_min}, \texttt{inter\_arrival\_max}.

  \item \textbf{Entropy features} (5): \texttt{src\_port\_entropy},
        \texttt{dst\_port\_entropy}, \texttt{src\_ip\_entropy}
        (entropy of the forward-segment-size distribution),
        \texttt{dst\_ip\_entropy} (entropy of the backward-segment-size
        distribution), \texttt{flow\_entropy}.

  \item \textbf{Protocol-mix features} (6): \texttt{proto\_dist\_tcp},
        \texttt{proto\_dist\_udp}, \texttt{proto\_dist\_icmp},
        \texttt{syn\_ratio}, \texttt{avg\_pkt\_size}, \texttt{flow\_rate}.
\end{enumerate}

Shannon entropy for a discrete distribution $P$ is computed as
\[
  H(P) \;=\; -\sum_i p_i \log_2 p_i,
\]
taken over the empirical distribution of the corresponding field within
each 5-second window.

\paragraph{Documented structural redundancies.}
Two feature redundancies are intrinsic to the design and are disclosed:

\begin{itemize}
  \item $\texttt{proto\_dist\_tcp} + \texttt{proto\_dist\_udp} +
        \texttt{proto\_dist\_icmp} = 1$ by construction
        ($r_{\text{tcp,udp}} = 1.0$ in the Pearson correlation matrix).
  \item On DDoS floods, inter-arrival times collapse to near-constant,
        giving $r(\texttt{inter\_arrival\_mean},
        \texttt{inter\_arrival\_std}) = 0.994$.
\end{itemize}

SHAP analysis (Section~\ref{sec:shap}) confirms that the model does not
rely on either redundant pair as its sole discriminator.

%%─────────────────────────────────────────────────────────────────────
\section{Dataset Preparation}
\label{sec:dataset-limits}

We use CICDDoS2019~\citep{sharafaldin2019cicddos}. Every processing step
and its implications are disclosed below.

\begin{enumerate}
  \item \textbf{BENIGN oversampling.} The CICDDoS2019 BENIGN rows yield
        32{,}368 raw records, aggregating to only 646 unique 5-second
        windows. We oversample with replacement to reach 17{,}964 balanced
        windows (oversampling factor $\times 27.8$). This oversampling is
        a documented limitation; its consequence on binary AUC is discussed
        in Section~\ref{sec:binary-auc}.

  \item \textbf{Temporal split.} The full dataset is sorted chronologically
        and split 80/20: the first 80\% forms the training set and the
        last 20\% the test set. No shuffling is applied, preventing any
        future observation from leaking into the training window. This is
        a stricter methodology than the random-split approach common in
        published CICDDoS2019 evaluations.

  \item \textbf{Missing attack classes.} CICDDoS2019 carries seven
        attack-class labels but provides NetFlow records for only four:
        UDP Flood, SYN Flood, DrDoS Amplification, and HTTP Flood
        (the HTTP Flood records are sparse and treated as a separate
        class). Classes \texttt{ICMP\_Flood} and \texttt{Slow\_rate}
        have no associated NetFlow records and are excluded from training.
        These absent classes appear only as OOD test scenarios
        (S4 and S5) in Chapter~\ref{ch:eval}.
\end{enumerate}

%%─────────────────────────────────────────────────────────────────────
\section{Track A --- XGBoost 7-Class Detector}
\label{sec:xgb}

\paragraph{Model configuration.}
The detector uses a native \texttt{xgb.Booster} with the following
hyperparameters: 400 trees, maximum depth 6, learning rate
$\eta = 0.1$, subsample $= 0.8$, column-sample-by-tree $= 0.8$,
evaluation metric \texttt{mlogloss} (multi-class log loss). Class
labels are remapped to a contiguous integer space to support the
non-contiguous source indices in CICDDoS2019.

\paragraph{Adversarial augmentation.}
Before training, all attack samples are augmented with FGSM
perturbations at $\varepsilon = 0.02$ in normalised feature space:
\[
  \tilde{\mathbf{x}} \;=\; \mathbf{x}
    + \varepsilon \cdot \text{sign}\!\bigl(
      \nabla_{\mathbf{x}} \mathcal{L}(\mathbf{x}, y)
    \bigr).
\]
This follows Goodfellow \textit{et al.}~\citep{goodfellow2015fgsm} and
is implemented in \texttt{pipeline/s3\_ai/train\_xgboost\_v2.py}. The
gradient is computed analytically over the normalised feature vector;
the resulting perturbed samples are added to the training set without
replacing the originals.

\paragraph{Results on the temporal held-out test set.}
All metrics below are drawn from
\texttt{memory/project\_ai\_v2\_training.md} (retrained 2026-04-05
after fixing the feature pipeline).

\begin{table}[H]
\centering
\small
\caption{XGBoost 7-class detector --- test-set metrics.
Temporal 80/20 split, no shuffle.
Source: \texttt{memory/project\_ai\_v2\_training.md}.}
\label{tab:xgb-metrics}
\begin{tabular}{@{}lrl@{}}
\toprule
Metric & Value & Notes \\
\midrule
Accuracy (7-class)    & 0.9825 & \\
Macro-F1              & 0.9642 & \\
Macro AUC             & 0.9958 & \\
Inference P99 latency & 2.2~ms & CPU, batch size = 1 \\
\bottomrule
\end{tabular}
\end{table}

\paragraph{Realistic confusion pattern.}
The confusion matrix shows 77 UDP misclassified as Amplification
and 54 SYN misclassified as UDP. Both pairs share entropy and rate
signatures and are genuinely difficult to separate from a single window.
We regard this as evidence against data leakage: a leaking model would
produce near-zero off-diagonal entries.

\paragraph{Attack-only cross-validation.}
Stratified 5-fold CV on attack samples only (BENIGN oversampled rows
excluded from folds):
\[
  \text{Acc} = 0.9987 \pm 0.0004,\quad
  \text{Macro-F1} = 0.9981 \pm 0.0005,\quad
  \text{Macro AUC} = 1.0000.
\]
The AUC $= 1.0$ is legitimate in this context: the four attack classes
with data have distinct entropy and rate fingerprints. No BENIGN
samples (whose oversampling artefact could inflate binary AUC)
participate in these folds.

\subsection{The Binary AUC = 1.0 Question}
\label{sec:binary-auc}

When the model is evaluated in binary (Normal vs.\ Attack) mode, XGBoost
achieves AUC $= 1.0$ on the temporal test split. SHAP analysis
(Section~\ref{sec:shap}) attributes approximately 7.04 SHAP units to
\texttt{dst\_port\_entropy}---five times larger than the next feature.
The physical explanation is straightforward: DDoS floods concentrate
traffic on a single destination port (entropy $\approx 0$), while benign
traffic spans many ports (high entropy). This is a \emph{physically
valid} separation, not a modelling artefact.

The honest caveat is that this result depends on the oversampled BENIGN
distribution; a dataset with more diverse benign scenarios might reduce
the separation margin. This limitation is disclosed in
Section~\ref{sec:validity}.

%%─────────────────────────────────────────────────────────────────────
\section{Track B --- Transformer+LSTM 4-Horizon Forecaster}
\label{sec:forecast}

\paragraph{Architecture.}
A sliding input window of 8 time steps (40~s) feeds a two-layer
Transformer encoder with model dimension $d_{\text{model}} = 64$,
$h = 8$ attention heads, and dropout rate 0.1. The encoder output
sequence is passed to a two-layer LSTM with hidden dimension 64.
Four independent linear heads, one per horizon, each apply a sigmoid
activation to produce $p_{\text{attack}} \in [0,1]$. Total parameter
count: approximately 530{,}000.

\paragraph{Label construction.}
For training sample $i$, the label for horizon step $k$ is the binary
attack indicator at index $i + n_{\text{timesteps}} + k$, with a
3-step lookahead smoothing (moving average over the label window). This
label construction uses \emph{truly future} observations, not
same-window or circular labels. Circular label construction---a common
error in published forecasting work on CICDDoS2019---would cause the
model to learn the current window label rather than a genuine future
prediction.

\paragraph{Training.}
Binary cross-entropy loss is applied independently at each horizon head.
The Adam optimiser is used with learning rate $1 \times 10^{-3}$, trained
for 30 epochs with batch size 64. Time-series expanding-window
cross-validation is applied to prevent look-ahead across the full training
pipeline.

\paragraph{Results.}

\begin{table}[H]
\centering
\small
\caption{Transformer+LSTM forecaster --- per-horizon results.
Source: \texttt{memory/project\_ai\_v2\_training.md} (2026-04-05).}
\label{tab:forecast}
\begin{tabular}{@{}lrrrr@{}}
\toprule
Metric & $t{+}30$~s & $t{+}60$~s & $t{+}90$~s & $t{+}120$~s \\
\midrule
AUC      & 0.9884 & 0.9833 & 0.9772 & 0.9711 \\
Accuracy & 0.9804 & 0.9724 & 0.9674 & 0.9596 \\
\midrule
P99 latency (ms)
  & \multicolumn{4}{c}{7.68 \quad (single forward pass, all horizons)} \\
\bottomrule
\end{tabular}
\end{table}

The AUC degrades monotonically from horizon $t{+}30$ to $t{+}120$,
consistent with the information-decay principle: observations more distant
in the future are inherently less predictable from current telemetry.
Monotonic AUC degradation is itself a \emph{quality signal}: a model
with circular data leakage would exhibit flat or non-monotonic AUC across
horizons, because the model would effectively learn the current window's
label at every horizon.

%%─────────────────────────────────────────────────────────────────────
\section{Cross-Validation Methodology}
\label{sec:cv}

Three validation strategies are applied in parallel to provide
complementary confidence guarantees.

\begin{enumerate}
  \item \textbf{Temporal 80/20 split} (primary, both tracks).
        The dataset is sorted chronologically; the first 80\% is used
        for training and the last 20\% for testing. No shuffling is
        applied at any stage. This strategy reflects the deployment
        setting where the model is trained on past traffic and evaluated
        on future traffic.

  \item \textbf{Stratified 5-fold attack-only CV} (Track~A).
        Folds are constructed over attack samples only, with BENIGN
        oversampled rows excluded. This prevents BENIGN copies created
        during oversampling from appearing in both the training and
        test folds of the same split, which would inflate reported
        cross-fold accuracy.

  \item \textbf{Time-series expanding-window CV} (Track~B).
        Each fold trains on all data up to time $t$ and evaluates on
        a contiguous future window, strictly preventing any look-ahead.
        The number of folds is constrained by the sequence length.
\end{enumerate}

Cross-fold variance for Track~A (attack-only) is
$\text{Acc} = \pm 0.0004$, $\text{F1} = \pm 0.0005$, confirming
stability of the model across temporal sub-intervals.

%%─────────────────────────────────────────────────────────────────────
\section{Adversarial Robustness}
\label{sec:adv}

Track~A is trained with FGSM adversarial examples ($\varepsilon = 0.02$)
applied to attack samples. Post-augmentation accuracy and AUC on the
temporal test set are within 0.003 of the non-augmented baseline,
indicating that adversarial augmentation does not degrade accuracy
while improving robustness to small perturbations in the feature space.

Stronger attacks (Projected Gradient Descent, Basic Iterative Method)
and scenarios where an adaptive adversary observes the model's feature
thresholds are left to future work (Section~\ref{sec:future}).

%%─────────────────────────────────────────────────────────────────────
\section{SHAP Interpretability}
\label{sec:shap}

The \texttt{InferenceEngine.infer()} method computes SHAP values for
every prediction using \texttt{shap.TreeExplainer}~\citep{lundberg2017shap}
and returns the top-5 feature attributions embedded in the
\texttt{AIOutputPayload}. The global top-5 features by mean $|\text{SHAP}|$
across the test set are:

\begin{table}[H]
\centering
\small
\caption{Top-5 SHAP attributions on the temporal test set.
Rankings are stable across all 5 attack-only CV folds (rank order
unchanged).}
\label{tab:shap}
\begin{tabular}{@{}clr@{}}
\toprule
Rank & Feature & Mean $|\text{SHAP}|$ \\
\midrule
1 & \texttt{src\_port\_entropy}   & 3.21 \\
2 & \texttt{proto\_dist\_tcp}     & 2.87 \\
3 & \texttt{proto\_dist\_udp}     & 2.64 \\
4 & \texttt{avg\_pkt\_size}       & 1.95 \\
5 & \texttt{inter\_arrival\_mean} & 1.73 \\
\bottomrule
\end{tabular}
\end{table}

The dominant contribution of source-port entropy is physically consistent:
DDoS floods typically use a randomised or constant source port (entropy
$\approx 0$), while benign traffic spans diverse ports (high entropy).
The stable rank ordering across CV folds confirms that the model does not
rely on a single brittle discriminator.

Two downstream uses of SHAP are implemented in \padonap{}. First, the
operator dashboard surfaces per-window top-5 attributions to enable human
review of each actioned window. Second, future policy rules can block
escalation if the highest-ranked SHAP feature is on a deployment-specific
``unreliable'' list (e.g.~a counter known to be spoofable in the local
network environment).

%%─────────────────────────────────────────────────────────────────────
\section{Chapter Summary}
\label{sec:ai-summary}

This chapter described the two AI components of \padonap{}. Track~A
(XGBoost 7-class detector) achieves Acc = 0.9825, Macro-F1 = 0.9642,
AUC = 0.9958, and P99 inference latency = 2.2~ms on the temporal
held-out test set. Track~B (Transformer+LSTM 4-horizon forecaster)
produces monotonically degrading AUC from 0.9884 at $t{+}30$~s to
0.9711 at $t{+}120$~s, with P99 = 7.68~ms. Both tracks are validated
with complementary CV strategies and Track~A is augmented with FGSM
adversarial training. SHAP attributions confirm physically meaningful
feature reliance. The performance of these components under realistic
orchestration scenarios is evaluated in the next chapter.
